\section{Introduction}\label{sec:introduction}

The entries of a matrix's first column are a row apart on the input tape
and consecutive in the output. A transposer may store coded information
as it works, and its head movements may depend on what it has read. If
we fix those movements, we can trace the possible dependencies of an
output bit. Fixing the movements also restricts the input. The question
is how much information survives the restriction.

For a fixed tape trajectory, a depth-reduction theorem supplies a small
set of steps whose values can be revealed so that the remaining
dependency paths have bounded length \cite[Corollary~O]{PaperFourteen}.
On the inputs sharing both the trajectory and the revealed values, each
output can then be recovered from a bounded set of input coordinates.
The recovered coordinate and its input copy can be far apart in the
incidence graph of that relation. Such separation is compatible with
binary storage codes of rate above one half
\cite[Corollary~A$'$]{PaperThirteen}. A linear-time transposer would
require rates arbitrarily close to one.

\Cref{thm:main} gives the degree dependence of this requirement. For
every sufficiently large degree cap it produces a family with rates
bounded below while its minimum paired distance tends to infinity. The
coordinates are bits of packets. Grouping bits this way gives our
strongest rate estimate. A coding upper bound with a larger deficit would
rule out the assumed machine. \Cref{cor:gap} makes the comparison precise,
including the success probability when the machine is correct only on
some inputs.

A second question concerns the individual bit positions of a physical
machine. \Cref{thm:simulator} groups source transitions into batches but
executes each head move and each append on an ordinary bit machine.
\Cref{thm:native-tail} uses its actual append times. The packet graph connects every bit of a permitted helper packet to
every bit of its target packet. The native graph follows each physical
append backwards through the simulator. On a set $G$ of correctly
transposed inputs, the native theorem bounds the input mass in fibres
whose rate falls appreciably below $\log_2|G|/L$. Each fibre keeps nearly
all of its trace's original bit positions. The native exponent is half
the packet exponent.

Finally, \cref{sec:patches} asks how much a local set of helpers can tell
us about the bit it is meant to recover. The answer accounts separately
for a bit's own bias and the information supplied by its helpers. A
fixed advantage over chance on a positive fraction of separated targets
requires increasingly large patch overlap as the code rate approaches
one. The outstanding question is whether a linear-time tape computation
would have to provide such an advantage at smaller overlap.

\paragraph{Relation to earlier work.}
The finite-control simulation of Hartmanis and Stearns evaluates
consecutive tape operations from a bounded neighbourhood
\cite[Theorem~2]{HartmanisStearns}. Hopcroft, Paul and Valiant synchronize
tape-block crossings \cite[\S2]{HPV1977}. Williams describes the
computation graph of an aligned machine using a constant number of
movement bits per block \cite[\S\S2--3]{Williams}. Here the sequence of batch descriptors determines the in-batch append
flags and when the terminal drain starts. The common output length
then fixes the number of bits appended during the drain. Earlier in this series, an external-memory
simulation fixed an event word at the scale of entire rows
\cite[Theorem~EM]{PaperEleven}. Our batch size is fixed before the
matrix dimension grows. Conditional entropy as a measure of imperfect
local recovery appears in earlier work on recoverable systems
\cite[\S IV]{ElishcoBarg}. Our prediction bound uses the separated-port
argument of \cite[Lemma~4.1]{PaperFive} and the soft-recovery inequality
for nested elimination \cite[Theorem~S]{SoftRecovery}.

\paragraph{Formal proofs.}
The Lean declarations named with each numbered result are its proofs.
The English outlines explain the constructions, estimates and choices of
parameters. \Cref{sec:verification} gives the verification procedure. The declaration index identifies the corresponding source modules.
